\documentclass[11pt,a4paper]{article}

\usepackage[utf8]{inputenc}
\usepackage[T1]{fontenc}
\usepackage[polish]{babel}
\usepackage{geometry}
\geometry{margin=2.5cm}

\usepackage{amsmath}
\usepackage{graphicx}
\usepackage{float}
\usepackage{booktabs}
\usepackage{caption}
\usepackage{pgfplots}
\pgfplotsset{compat=1.18}
\usepackage{listings}
\usepackage{xcolor}
\usepackage{parskip}
\usepackage{enumitem}


\lstset{
    language=C++,
    basicstyle=\ttfamily\small,
    backgroundcolor=\color{white},
    frame=single,
    rulecolor=\color{black},
    breaklines=true,
    tabsize=2,
    numbers=left,
    numberstyle=\tiny\color{gray},
    captionpos=b,
    keepspaces=true,
    float=hbtp,
    aboveskip=8pt,
    belowskip=8pt
}

\title{\textbf{Analiza algorytmów}\\
\vspace{0.3cm}
\large Maciej Pluta, 287339}
\date{}

\begin{document}
\maketitle


\section{Wprowadzenie}
Projekt obejmował implementację i szczegółową analizę czterech fundamentalnych algorytmów. Każdy algorytm został zaimplementowany w kilku wariantach, co pozwoliło na porównanie ich wydajności, poprawności oraz praktycznej użyteczności w różnych scenariuszach.

Przeanalizowane algorytmy:
\begin{itemize}[itemsep=2pt, topsep=2pt]
    \item \textbf{CUT-ROD} – algorytm optymalizacji cięcia pręta wykorzystujący programowanie dynamiczne
    \item \textbf{LCS (Longest Common Subsequence)} – algorytm znajdowania najdłuższego wspólnego podciągu
    \item \textbf{ACTIVITY SELECTOR} – algorytm wyboru maksymalnego zbioru niekonfliktujących zajęć
    \item \textbf{HUFFMAN} – algorytm konstrukcji optymalnego kodu prefiksowego (wersja ternarna)
\end{itemize}


\section{Kluczowe fragmenty implementacji}

\subsection{CUT-ROD}
\textbf Algorytm optymalizacji cięcia pręta na kawałki w celu maksymalizacji zysku. Mając pręt długości $n$ i tablicę cen $p[i]$ za pręt długości $i$, algorytm znajduje optymalny podział. Kluczową cechą jest optymalna podstruktura.

\begin{lstlisting}[caption={CUT-ROD – wersja Bottom-Up}]
int cutRodBottomUp(const vector<int>& p, int n, vector<int>& s) {
    vector<int> dp(n+1, 0);
    for (int j = 1; j <= n; j++) {
        int best = INT_MIN;
        for (int i = 1; i <= j; i++) {
            if (p[i-1] + dp[j-i] > best) {
                best = p[i-1] + dp[j-i];
                s[j] = i;
            }
        }
        dp[j] = best;
    }
    return dp[n];
}
\end{lstlisting}

\textbf{Działanie:} Algorytm wypełnia tablicę \texttt{dp} od najmniejszych do największych długości. Dla każdej długości $j$ sprawdza wszystkie możliwe pierwsze cięcia $i$, wybierając to, które daje maksymalny zysk. 

\textbf{Złożoność:} $O(n^2)$ czasowa, $O(n)$ pamięciowa

\subsection{LCS}
\textbf Algorytm znajduje najdłuższy podciąg wspólny dla dwóch sekwencji przy wykorzystaniu tablicy dwuwymiarowej. Zaimplementowano dwie wersje: rekurencyjną ze spamiętywaniem i iteracyjną.

\begin{lstlisting}[caption={LCS – wersja iteracyjna}]
string lcsIter(const string& a, const string& b) {
    int n = a.size(), m = b.size();
    vector<vector<int>> dp(n+1, vector<int>(m+1,0));
    
    for (int i = 1; i <= n; i++)
        for (int j = 1; j <= m; j++)
            if (a[i-1] == b[j-1])
                dp[i][j] = dp[i-1][j-1] + 1;
            else
                dp[i][j] = max(dp[i-1][j], dp[i][j-1]);
    
    string res;
    int i = n, j = m;
    while (i > 0 && j > 0) {
        if (a[i-1] == b[j-1]) {
            res = a[i-1] + res;
            i--; j--;
        } else if (dp[i-1][j] >= dp[i][j-1]) i--;
        else j--;
    }
    return res;
}
\end{lstlisting}

\textbf{Działanie:} Tablica \texttt{dp[i][j]} przechowuje długość LCS dla prefiksów. Jeśli znaki są równe, długość zwiększa się o 1. W przeciwnym razie wybierane jest lepsze z rozwiązań dla krótszych prefiksów.

\textbf{Złożoność:} $O(n \cdot m)$ czasowa, $O(n \cdot m)$ pamięciowa

\subsection{ACTIVITY SELECTOR}
\textbf Problem wyboru zajęć polega na znalezieniu maksymalnego zbioru niekonfliktujących aktywności. Algorytm zachłanny sortuje zajęcia po czasie zakończenia i wybiera kolejne dostępne zajęcia. Strategia ta gwarantuje optymalność dzięki własności zachłannego wyboru.

\begin{lstlisting}[caption={ACTIVITY SELECTOR – wersja zachłanna}]
vector<pair<int,int>> activityIterative(vector<pair<int,int>> a) {
    sort(a.begin(), a.end(), [](auto& x, auto& y) {
        return x.second < y.second;
    });

    vector<pair<int,int>> res;
    int last = 0;
    for (auto& x : a) {
        if (x.first >= last) {
            res.push_back(x);
            last = x.second;
        }
    }
    return res;
}
\end{lstlisting}

\textbf{Działanie:} Algorytm sortuje zajęcia rosnąco po czasie zakończenia. Następnie iteracyjnie wybiera każde zajęcie, które zaczyna się po zakończeniu poprzednio wybranego. Zawsze wybiera zajęcie kończące się najwcześniej spośród dostępnych.

\textbf{Złożoność:} $O(n \log n)$ czasowa (dominowane przez sortowanie), $O(n)$ pamięciowa

\subsection{HUFFMAN}
\textbf Algorytm konstruuje optymalny kod prefiksowy minimalizujący średnią długość zakodowanej wiadomości. Wersja ternarna używa trzech symboli (0,1,2) zamiast dwóch standardowych. Algorytm buduje drzewa łączenia węzłów o najmniejszych częstotliwościach.

\begin{lstlisting}[caption={HUFFMAN – wersja ternarna}]
map<char,string> huffman(const map<char,long long>& freq) {
    priority_queue<Node*, vector<Node*>, Cmp> pq;
    for (auto& x : freq)
        pq.push(new Node(x.second, x.first));

    const int k = 3;
    while (pq.size() > 1) {
        int take = min(k, (int)pq.size());
        Node* n = new Node(0);
        for (int i = 0; i < take; i++) {
            auto t = pq.top(); pq.pop();
            n->ch.push_back(t);
            n->f += t->f;
        }
        pq.push(n);
    }

    map<char,string> res;
    Node* root = pq.top();
    gen(root, "", res);
    del(root);
    return res;
}
\end{lstlisting}

\textbf{Działanie:} Kolejka priorytetowa przechowuje węzły posortowane rosnąco po częstotliwościach. W każdej iteracji łączone są 3 węzły o najmniejszych częstotliwościach. Nowy węzeł otrzymuje sumę częstotliwości dzieci. Proces kończy się gdy w kolejce pozostanie tylko korzeń drzewa.

\textbf{Złożoność:} $O(n \log n)$ czasowa, $O(n)$ pamięciowa


\section{Wyniki testów i analiza}

\subsection{Testy poprawności}

\begin{table}[H]
\centering
\caption{Wyniki testów poprawności}
\begin{tabular}{ll}
\toprule
\textbf{Algorytm / Test} & \textbf{Wynik} \\
\midrule
CUT-ROD (n=4) & Naive: 10, Memo: 10 | 2 2, Bottom-Up: 10 | 2 2 \\
LCS & "BCBA" (obie implementacje) \\
ACTIVITY SELECTOR & (1,4) (5,7) (8,11) \\
HUFFMAN & a: 120, b: 121, c: 122, d: 10, e: 11, f: 0 \\
\bottomrule
\end{tabular}
\end{table}

\textbf{Komentarz:} Wszystkie implementacje każdego algorytmu zwracają identyczne wyniki, co potwierdza ich poprawność. Dla CUT-ROD wszystkie trzy wersje zwracają maksymalny zysk 10, a wersje Memo i Bottom-Up dodatkowo odtwarzają optymalny podział 2+2.

\subsection{Testy wydajnościowe}

\begin{table}[H]
\centering
\caption{Czas wykonania w mikrosekundach [µs]}
\begin{tabular}{llll}
\toprule
\textbf{Algorytm} & \textbf{Wersja 1} & \textbf{Wersja 2} & \textbf{Komentarz} \\
\midrule
CUT-ROD (n=20) & Memo: 4 & Bottom-Up: 6 & Wersja rekurencyjna ze spamiętywaniem jest szybsza \\
CUT-ROD (n=30) & Memo: 5 & Bottom-Up: 3 & Wersja iteracyjna zaczyna być szybsza \\
CUT-ROD (n=40) & Memo: 9 & Bottom-Up: 5 & Różnica na korzyść Bottom-Up staje się wyraźna \\
ACTIVITY (2000) & Greedy: 481 & DP: 31554 & Algorytm zachłanny jest 65× szybszy \\
LCS (500x500) & Iteracyjny: 3000 & --- & Dobry czas dla dużych danych \\
\bottomrule
\end{tabular}
\end{table}

\textbf{Komentarz ogólny:} Testy wydajnościowe pokazują wyraźną przewagę wersji iteracyjnych i algorytmów zachłannych dla dużych zbiorów danych. Algorytm naiwny CUT-ROD nie był testowany dla n>10 ze względu na złożoność wykładniczą.


\section{Wizualizacja wyników}

\begin{figure}[H]
\centering
\begin{tikzpicture}
\begin{axis}[
    width=0.8\textwidth, height=0.4\textwidth,
    xlabel={Liczba aktywności (n)}, ylabel={Czas wykonania [µs]},
    legend pos=north west, grid=major,
    title={ACTIVITY SELECTOR: Greedy vs DP},
    ymode=log,
    log basis y=10,
    ymin=10, ymax=100000
]
\addplot[blue, mark=*, thick] coordinates {(100,50) (500,150) (1000,280) (2000,481)};
\addlegendentry{Greedy $O(n \log n)$}

\addplot[red, mark=square*, thick] coordinates {(100,200) (500,2500) (1000,9500) (2000,31554)};
\addlegendentry{DP $O(n^2)$}
\end{axis}
\end{tikzpicture}
\end{figure}

\begin{figure}[H]
\centering
\begin{tikzpicture}
\begin{axis}[
    width=0.8\textwidth, height=0.4\textwidth,
    xlabel={Długość pręta (n)}, ylabel={Czas [µs]},
    legend pos=north west, grid=major,
    title={CUT-ROD: Memo vs Bottom-Up},
    ymin=0
]
\addplot[blue, mark=*, thick] coordinates {(20,4) (30,5) (40,9)};
\addlegendentry{Memo (rekurencyjna)}

\addplot[red, mark=square*, thick] coordinates {(20,6) (30,3) (40,5)};
\addlegendentry{Bottom-Up (iteracyjna)}
\end{axis}
\end{tikzpicture}
\end{figure}


\section{Wnioski}

\subsection*{CUT-ROD}
\textbf{Zalety:} Gwarantuje rozwiązanie optymalne, prosty w implementacji, łatwe odtwarzanie rozwiązania.\\
\textbf{Wady:} Złożoność kwadratowa $O(n^2)$, wersja naiwna ma złożoność wykładniczą.\\
\textbf{Zastosowanie kiedy:} Dla małych i średnich n, gdy potrzebne jest dokładne optymalne rozwiązanie.\\
\textbf{Wersja preferowana:} Bottom-Up dla n > 25, Memo dla mniejszych danych.

\subsection*{LCS}
\textbf{Zalety:} Gwarantuje optymalne rozwiązanie, dobrze skalowalny dla umiarkowanych danych.\\
\textbf{Wady:} Wysokie zużycie pamięci $O(n \cdot m)$, wolniejszy dla bardzo długich sekwencji.\\
\textbf{Zastosowanie kiedy:} Porównywanie sekwencji tekstowych lub genetycznych o umiarkowanej długości.\\
\textbf{Wersja preferowana:} Iteracyjna - bardziej stabilna wydajnościowo.

\subsection*{ACTIVITY SELECTOR}
\textbf{Zalety:} Algorytm zachłanny jest bardzo szybki $O(n \log n)$, prosty w implementacji.\\
\textbf{Wady:} DP ma złożoność $O(n^2)$, co czyni go niepraktycznym dla dużych danych.\\
\textbf{Zastosowanie kiedy:} Zawsze preferować wersję zachłanną - jest optymalna i znacznie szybsza.\\
\textbf{Różnica wydajności:} Greedy 65× szybszy niż DP dla 2000 zajęć.

\subsection*{HUFFMAN}
\textbf{Zalety:} Generuje optymalne kody prefiksowe, wersja ternarna działa poprawnie.\\
\textbf{Wady:} Konieczność przechowywania drzewa kodowego, koszt kodowania/odkodowania.\\
\textbf{Zastosowanie kiedy:} Kompresja danych, gdy częstotliwości symboli są znane i nierównomierne.

\subsection*{Ogólne wnioski}
\begin{enumerate}[itemsep=2pt, topsep=2pt]
    \item \textbf{Iteracyjne vs rekurencyjne:} Wersje iteracyjne (Bottom-Up) są szybsze dla dużych danych dzięki brakowi narzutu rekurencyjnego.
    
    \item \textbf{Zachłanne vs dynamiczne:} Algorytmy zachłanne są znacznie szybsze gdy problem ma własność zachłannego wyboru (ACTIVITY SELECTOR 65× szybszy).
    
    \item \textbf{Skalowalność:} CUT-ROD $O(n^2)$ słabo skalowalny, LCS $O(n \cdot m)$ umiarkowanie, Greedy ACTIVITY $O(n \log n)$ dobrze.
    
    \item \textbf{Praktyczność:} Wersje naiwne (rekurencyjne bez optymalizacji) są niepraktyczne nawet dla średnich danych.
    
    \item \textbf{Wybór algorytmu:} Zależy od rozmiaru danych, wymagań wydajnościowych i struktury problemu.
\end{enumerate}

\textbf{Podsumowanie:} Projekt potwierdził, że nie ma uniwersalnie najlepszego algorytmu. Wybór między programowaniem dynamicznym a podejściem zachłannym zależy od specyfiki problemu. Dla problemów z optymalną podstrukturą i nakładającymi się podproblemami (CUT-ROD, LCS) DP jest konieczne. Dla problemów z własnością zachłannego wyboru (ACTIVITY SELECTOR, HUFFMAN) algorytmy zachłanne są znacznie efektywniejsze.

\end{document}